\ExerciseSection

¶ 1) a) 设 \((x_{k})_{1\leqslant k\leqslant m}\) 是一个实数的有限序列，满足 \(|x_k|\leqslant 1\) （对 \(1\leqslant k\leqslant m\) ），且 \(\sum_{k = 1}^{m}x_{k} = 0\) 。证明存在一个置换 \(\sigma\) ，即区间 \([1,m]\) （属于 \(\mathbf{N}\) ）的一个置换，使得

\[
\left| \sum_ {k = 1} ^ {p} x _ {\sigma (k)} \right| \leqslant 1
\]

对一切 \(p = 1, 2, \ldots, m\) 成立，并且该置换保持诸指标 \(h\) （对它们 \(x_h > 0\) ）的次序以及诸指标 \(h\) （对它们 \(x_h < 0\) ）的次序{[}即，若 \(h < k\) 且 \(x_h > 0\) 且 \(x_k > 0\) ，则 \(\sigma(h) < \sigma(k)\) ，对指标 \(h\) 满足 \(x_h < 0\) 者亦然{]}。

\begin{enumerate}
\def\labelenumi{\alph{enumi})}
\setcounter{enumi}{1}
\tightlist
\item
  设 \((\pmb{x}_k)_{1\leqslant k\leqslant m}\) 是点 \(\pmb{x}_k = (x_{ki})_{1\leqslant i\leqslant n}\) （属于 \(\mathbf{R}^n\) ）的一个有限序列，满足 \(||\pmb{x}_k||\leqslant 1\) （对 \(1\leqslant k\leqslant m,\) ），且 \(\sum_{k = 1}^{m}\pmb{x}_k = 0.\) 证明存在一个置换 \(\sigma\) ，即区间 \([1, m]\) （属于 \(\mathbf{N}\) ）的一个置换，使得
\end{enumerate}

\[
\left\| \sum_ {k = 1} ^ {p} x _ {\sigma (k)} \right\| \leqslant 5 ^ {(n - 1) / 2} \text { for   all } p = 1, 2, \dots , m.
\]

{[}对 \(n\) 用归纳法证明，把 \(\mathbf{R}^n\) 看作乘积 \(\mathbf{R}^{n-1} \times \mathbf{R}\) ，并令 \(x_k = (x_k', x_{kn})\) ，其中 \(x_k' \in \mathbf{R}^{n-1}\) 。取一个子集 \(\mathbf{H}\) ，即区间 {[}1, m{]} （属于 N ）的一个子集，使 \(\left\|\sum_{k\in H}x_{k}\right\|\) 达到最大；借助一个旋转，化归为 \(\sum_{k\in H}x_{k}'=0\) 的情形，并证明此时必有 \(x_{kn}\geqslant0\) （对 \(k\in H\) ）且 \(x_{kn}\leqslant0\) （对 \(k\notin H\) ）；然后利用 a) 与归纳假设。{]}

\begin{enumerate}
\def\labelenumi{\alph{enumi})}
\setcounter{enumi}{2}
\tightlist
\item
  设 \((\pmb{x}_k)_{1 \leqslant k \leqslant m}\) 是 \(\mathbf{R}^n\) 的点的一个有限序列，满足 \(||\pmb{x}_k|| \leqslant 1\) （对 \(1 \leqslant k \leqslant m\) ），且 \(\left\| \sum_{k=1}^{m} \pmb{x}_k \right\| = a > 0\) 。证明存在一个置换 \(\sigma\) ，即区间 \([1, m]\) （属于 \(\mathbf{N}\) ）的一个置换，使得
\end{enumerate}

\[
\left\| \sum_ {k = 1} ^ {p} x _ {\sigma (k)} \right\| \leqslant (a + 1) 5 ^ {(n - 1) / 2}
\]

对一切 \(p\) （满足 \(1 \leqslant p \leqslant m\) ）成立{[}化归为情形 \(b\)。{]}

\begin{enumerate}
\def\labelenumi{\arabic{enumi})}
\setcounter{enumi}{1}
\tightlist
\item
  设 \((\pmb{x}_m)_{m\in \mathbb{N}}\) 是 \(\mathbf{R}^n\) 的点的一个无穷序列，满足 \(\lim_{m\to \infty}\pmb{x}_m = 0\) 。对每个有限子集 \(\mathbf{H}\) （\(\mathbf{N}\) 的），令 \(S_{\mathbf{H}} = \sum_{m\in \mathbf{H}}x_m\) 。证明只有两种可能：或者
\end{enumerate}

\begin{enumerate}
\def\labelenumi{(\roman{enumi})}
\item
  \(\lim_{\mathfrak{F}}||\mathbf{S}_{\mathbf{H}}|| = +\infty\) ，其中 \(\mathfrak{F}\) 是 \(\mathbf{N}\) 的所有有限子集组成的有向集；或者：
\item
  存在置换 \(\sigma\) （\(\mathbf{N}\) 的），使得通项为 \(x_{\sigma(m)}\) 的级数在 \(\mathbf{R}^n\) 中收敛；此时，集 \(\mathbf{A}\) 由这些级数的和 \(\sum_{m=0}^{\infty} x_{\sigma(m)}\) （对所有具有此性质的置换 \(\sigma\) ）组成，是 \(R^n\) 中的一个仿射线性簇。{[}首先借助习题 1 b) 证明，每个（关于 \(\mathfrak{F}\) 的）映射 \(H \to S_H\) 的聚点都是某个收敛级数 \((x_{\sigma(m)})\) 对适当置换 \(\sigma\) 的和。利用第三章 § 5 习题 3 证明：集 A 若非空，则是 \(R^n\) 中某个闭子群的一个陪集。最后借助习题 1 c) 与第二章 § 4 习题 15 证明 A 是连通的。{]}
\end{enumerate}

